% package
\usepackage{alltt}
\usepackage{amsfonts}
\usepackage{amsmath}
\usepackage{hyphenat}
\usepackage{amssymb}
\usepackage{amsthm}
\usepackage{bm}
\usepackage{booktabs}
\usepackage{caption}
\usepackage{color}
% \usepackage[colorlinks]{hyperref}
\usepackage{enumitem}
\usepackage{etoolbox}
\usepackage{fancyhdr}
\usepackage{flafter}
\usepackage{geometry}
\usepackage{graphicx}
\usepackage{mathalfa}
\usepackage{mathdots}
\usepackage{mathrsfs}
\usepackage{mathtools}
% \usepackage{microtype}
\usepackage{multirow}
\usepackage{multicol}
\usepackage{pdflscape}
\usepackage{pgffor}
\usepackage{pgfplots}
\usepackage{siunitx}
\usepackage{slashed}
\usepackage{subfigure}
\usepackage{tabularx}
\usepackage{tikz}
\usepackage{tkz-euclide}
\usepackage[normalem]{ulem}
\usepackage{xcolor}
\usepackage{xparse}
\usepackage{xpatch}
\usepackage[all]{xy}
\usepackage{imakeidx}
\usepackage[most]{tcolorbox}

% --- Custom Color Definitions ---
% These closely match the green tones in your image
\definecolor{posterBorderGreen}{RGB}{91, 131, 80}      % Darker green for borders and text
\definecolor{posterHeaderBg}{RGB}{216, 235, 200}     % Light green background for the header

% --- Custom Environment Definition ---
\newtcolorbox{hyperbolicbox}[1]{
    enhanced,
    colback=white,                           % White body background
    colframe=posterBorderGreen,              % Outer container border color
    boxrule=3pt,                             % Border line thickness
    arc=6pt,                                 % Rounded corner radius
    boxsep=10pt,                             % Padding inside the main container
    top=5pt, bottom=10pt, left=10pt, right=10pt,
    halign=flush left,
    breakable,
    % This handles the embedded, full-width header block at the top
    before upper={
        \begin{tcolorbox}[
            colback=posterHeaderBg,
            colframe=posterBorderGreen,
            boxrule=1.5pt,
            arc=4pt,
            width=\linewidth,
            top=8pt, bottom=8pt, left=10pt, right=10pt,
        ]
        \centering{\color{posterBorderGreen}\sffamily\bfseries\large #1}
        \end{tcolorbox}
    }
}




\theoremstyle{definition}
\newtheorem{thm}{Theorem}
\newtheorem{defi}[thm]{Definition}

\newcommand{\C}{\mathbb{C}}
\newcommand{\R}{\mathbb{R}}
\newcommand{\Q}{\mathbb{Q}}
\newcommand{\Z}{\mathbb{Z}}
\newcommand{\N}{\mathbb{N}}
\renewcommand{\H}{\mathbb{H}}
\newcommand{\F}{\mathbb{F}}

\newcommand{\topo}{\mathcal{T}}
\newcommand{\B}{\mathcal{B}}
\renewcommand{\S}{\mathcal{S}}
% Symbols
\newcommand{\noteq}{\not=}
\newcommand{\re}{\mathrm{Re\,}}
\newcommand{\im}{\mathrm{Im\,}}
\newcommand{\nexist}{\not\!\exists\,}
\newcommand{\abs}[1]{\left\lvert #1\right\rvert}
\newcommand{\norm}[1]{\left\lVert #1\right\rVert}
\newcommand{\floor}[1]{\left\lfloor #1\right\rfloor}
\newcommand{\ceil}[1]{\left\lceil #1\right\rceil}
\newcommand{\set}[1]{\left\{#1\right\}}
\newcommand{\closure}[2][3]{%
{}\mkern#1mu\overline{\mkern-#1mu#2}}
\newcommand{\interior}{\mathrm{Int\,}}
\newcommand{\upint}{\mathchoice%
    {\mkern13mu\overline{\vphantom{\intop}\mkern7mu}\mkern-20mu}%
    {\mkern7mu\overline{\vphantom{\intop}\mkern7mu}\mkern-14mu}%
    {\mkern7mu\overline{\vphantom{\intop}\mkern7mu}\mkern-14mu}%
    {\mkern7mu\overline{\vphantom{\intop}\mkern7mu}\mkern-14mu}%
  \int}
\newcommand{\lowint}{\mkern3mu\underline{\vphantom{\intop}\mkern7mu}\mkern-10mu%
\int}
\newcommand{\aand}{\quad\text{and}\quad}
\newcommand{\res}[1]{\!\mid_{#1}}
\newcommand{\normal}{\unlhd}
\newcommand{\group}[1]{\langle#1\rangle}
\renewcommand{\tt}[1]{\quad\text{#1}\quad}
\renewcommand{\d}{\mathop{}\!\mathrm{d}} 
\renewcommand{\emptyset}{\varnothing}
\DeclareMathOperator{\lcm}{lcm}
\DeclareMathOperator{\diam}{diam}
\DeclareMathOperator{\id}{id}
\DeclareMathOperator{\Bd}{Bd}